% sec_identification.tex

\subsection{Additional Identification Evidence}
\label{sec:additional_id}

\paragraph{FL entry event study.}
Within-PBU event studies around first FL entry
(Figure~\ref{fig:fl_entry_es}, Appendix~\ref{app:robustness}) show
positive pre-event coefficients at $t = -2$ and $t = -3$: PBUs that
eventually receive FL entry \emph{already had higher prices}. Two
readings are possible.

The first is strategic market selection. Cartels pick competitive,
high-value markets---the pre-trend reflects that selection---and then
deploy cover bidders to sustain elevated prices. Under this reading,
the pre-trend is part of the story rather than a threat to it, and
the item--year--PBU fixed effects in Equation~\eqref{eq:ols_main}
absorb the price \emph{level} in each market.

The second is straightforward selection bias. If unobserved market
characteristics attract both FL firms and higher prices, the
fixed effects may not fully absorb time-varying within-market
confounders, and the OLS overstates the FL--price association.

We cannot adjudicate between the two readings, but two pieces of
evidence bear on both readings. The reverse-causality test
(Section~\ref{sec:mechanisms}, M3) shows that the elasticity of FL
entry with respect to lagged prices is tiny ($\approx 0.004$), which
limits the scope for price-chasing selection. And the sensitivity
analysis ($RV = 17.5\%$) implies that a confounder driving the
pre-trend would need to explain a large share of residual variation in
both treatment and outcome. Still, we regard the pre-trends as an
unresolved ambiguity---one more reason to interpret the estimates as
conditional associations rather than causal effects.

\paragraph{Minimum-bidder constraint variation.}
In convite, when $n_{\text{genuine}} < 3$ the cartel \emph{must}
deploy cover bidders (corner solution). When $n \geq 3$, cover bidding
is voluntary. Interacting FL with a constraint-binding indicator within
the convite subsample yields $\hat{\beta}_{\text{FL}}(n \geq 3) =
0.076$ (SE~$= 0.021$, $p < 0.001$)---a 7.6\% premium from
\emph{voluntary} cover bidding---and $\hat{\beta}_{\text{FL} \times
(n<3)} = -0.160$ ($p < 0.001$). The price effect concentrates in
voluntary deployments, not constraint-driven ones.

\paragraph{CADE enforcement shock.}
A Callaway--Sant'Anna DiD exploiting CADE conviction timing shows that
FL presence drops sharply after conviction (ATT~$= -0.275$,
$p < 0.001$), as one would expect if cover bidders exit when the
cartel is prosecuted. The price ATT is positive but imprecise (0.145,
SE~$= 0.110$).
